\documentclass[10pt]{exam}
\usepackage[phy]{template-for-exam}
\usepackage{my-tikz-clipart}
\usetikzlibrary{shadings,decorations.pathmorphing,arrows.meta,patterns}

\title{Projectiles \#4}
\author{Rohrbach}
\date{\today}

\begin{document}
\maketitle

\begin{questions}

\question
  An air rocket is launched with an initial $x$-velocity of 17 m/s and an initial $y$-velocity of 13 m/s.
  
  \begin{parts}
    \part
      Label the diagram of the situation and write down knowns and unknowns.

      \vspace{2em}
    
      \def\posscale{0.55}
      \def\vox{17*\posscale}
      \def\voy{13*\posscale}
      \begin{tikzpicture}[
        declare function ={
          trajx(\t) = \vox*\t;
          trajy(\t) = \voy*\t-4.9*\t*\t;
          vy(\t) = \voy-9.8*\t;
          th(\t) = atan(vy(\t)/(\vox));
        }
      ]

        \def\vectorscale{0.3}

        \draw[-Stealth,very thick, red] 
          (0,0) -- 
          ++({\vox*\vectorscale},{\voy*\vectorscale})
          coordinate (end of v);
        \draw[dashed,red] 
          (0,0) -| (end of v);
        \draw[dotted,thick] 
          plot[
            samples=100, 
            domain=0:(2*\voy/9.8), 
            smooth, 
            variable=\t
          ] 
            ({trajx(\t)}, {trajy(\t)});
        \draw[thick] (-0.2,-0.1) -- ++(14.2,0);

        \def\myt{0.6}

        \begin{scope}[
          shift={({trajx(\myt)}, {trajy(\myt)})},
          rotate={-90+th(\myt)}
        ]
          \draw[fill=gray] (-0.05,0.2) rectangle (0.05,-0.2);
          \draw[fill=gray!10] (-0.05,0.2) -- (0,0.35) -- (0.05,0.2);
        \end{scope}
      \end{tikzpicture}

      \vs
    
    \part
      Calculate the {\bf magnitude} and {\bf angle} of the rocket's initial resultant velocity.
      \vs
    
    \part
      Calculate the maximum height of the rocket.
      \vs
    
    \part
      Calculate the amount of time that the rocket is in the air.
      \vs
    
    \part
      What is the overall $x$-displacement of the rocket before it hits the ground?
      \vs

  \end{parts} 

\pagebreak
  

\question
  A Hollywood stunt woman drives a car horizontally off a cliff that is 94 m high.  With what velocity must she drive the car so that it lands on a barge located 250 m away from the bottom of the cliff (measured horizontally)?

  \begin{parts}
    \part 
      Label the diagram of the situation and write down knowns and unknowns.

      \vspace{2em}
    
      \begin{tikzpicture}
        \draw[fill=gray] (3,-0.2)
          -- ++(0,2.2)
          -- ++(-2,0)
          -- ++(0,-2.2)
          -- cycle;
        \draw[-Stealth,very thick, red] (3, 2.3) -- ++(1,0);
        \draw[dotted,thick] (3.1,2.3) parabola (10,0.3) ++ (0,-0.1) coordinate (end);

        \draw[ultra thin, gray] (2,2.6) -- ++(-.3,0);
        \draw[ultra thin, gray] (1.9,2.4) -- ++(-.3,0);
        \draw[ultra thin, gray] (1.95,2.2) -- ++(-.4,0);

        \draw (2,2.1) pic[scale=0.2] {car};

        \draw[
          thick,
          blue,
          rounded corners,
          decorate,
          decoration={
            random steps,
            segment length=2mm,
            amplitude=2pt
          }] 
          (3,0) -- (12.5,0);
        \draw[fill=gray, rounded corners] 
          (end) ++(-0.7,0) -- ++(1.4,0) 
          -- ++ (0.2,0.2) -- ++(0,-0.6) -- ++ (-1.8,0)
          -- ++ (0,0.6)
          -- cycle;

      \end{tikzpicture}

      \vspace{2em}
    \part
      How much time will it take for the car to fall down the cliff?
      \vs
  
    \part
      What velocity did the woman need to have going off the cliff?
      \vs
  
    \part
      What is the woman's $y$-velocity just before she hits the ground?
      \vs
  
    \part
      Draw a triangle to find the woman's {\bf resultant final velocity} and the {\bf angle}.
      \vs

    
    \uplevel{\footnotesize Answers: 
      (b) 4.38 sec; 
      (c) 57.1 m/s; 
      (d) 42.9 m/s;
      (e) 71.4 m/s @ $36.9^\circ$.
    }

  \end{parts}

\pagebreak

\end{questions}






\end{document}